\section{Three motivating facts}
\label{sec:facts.sec}

This section discusses three facts about subjective expectations that motivate our model. 

\paragraph{Fact 1. Expectations are volatile.}
As \citet{shiller1981stock} famously observed, stock dividends and earnings are too stable to explain the large swings in asset prices. In contrast, recent survey evidence shows that \textit{subjective expectations} of future cash flows are highly volatile, offering a potential explanation for the stock market's excess volatility \citep{bordalo2020expectations, de_la_o_subjective_2021}.

We illustrate the magnitude of the difference in the volatility of subjective and objective expectations by constructing an "objective" forecast---based on the statistical properties of dividends---against a survey-based "subjective" forecast. Specifically, we estimate an AR(1) process for quarterly aggregate dividend and earnings growth. The predicted one-year-ahead dividend growth from this AR(1) serves as our measure of objective expectations, $\mathbb{E}^{obj}[\Delta d_{t,t+4}]$. We then compute the fraction of the total variance in realized dividend growth expectations relative to the variance of realized dividend growth. We construct an analogous measure for subjective expectations, $\mathbb{E}^{sub}[\Delta d_{t,t+4}]$, by using the survey of analyst expectations from I/B/E/S. The differences in volatility are significant:
\begin{equation}
    \frac{\mathrm{Var}\!\bigl[\mathbb{E}^{obj}[\Delta d_{t,t+4}]\bigr]}{\mathrm{Var}[\Delta d_{t,t+4}]} = 0.05,
    \qquad\quad
    \frac{\mathrm{Var}\!\bigl[\mathbb{E}^{sub}[\Delta d_{t,t+4}]\bigr]}{\mathrm{Var}[\Delta d_{t,t+4}]} = 0.70.
    \nonumber
\end{equation}
Whereas objective expectations represent a negligible fraction of the total variation in dividend growth, subjective expectations are an order of magnitude more volatile.\footnote{The small fraction of variation explained by objective expectations is consistent with the estimated autocorrelation of cash flows in, for example, \citet{bansal2004risks}.}

% Next, we provide a sense of the importance of belief heterogeneity, we investigate the link between investor heterogeneity and  turnover volume. Turnover is informative about the importance of our mechanism: if investors forecast errors are random across agents and time, stock turnover should be zero and thus, a model with homogeneous beliefs should be a good representation of reality. In turn, if a set of investors is rational turnover should also be zero. If turnover and market returns are correlated, this is evidence of belief heterogeneity.

% This discussion motivates our model. With a model, we can then provide an explicit relation between belief heterogeneity and stock market turnover to better approximate the importance of belief heterogeneity in explaining asset prices.  

% we argued that there is some inconsistency in the volatility predicted by DLM. I forgot what it was, but I'm wondering if we can use some convex combination: $\mathbb{E}_{t}^{m}=\beta*\mathbb{E}_{t}^{i}+(1-\beta)*\mathbb{E}_{s}^{i}$ to be estimated (estimate beta) and then argue they matter.

\paragraph{Fact 2. Expectations are heterogeneous.}
Fact 1 provides evidence of deviations from full-information rational expectations (FIRE). 
Our business cycle model emphasizes how heterogeneity in these expectations leads to business cycle amplification through asset markets. 
Critical to our theory is the presence of heterogeneity in beliefs. 
Concrete evidence is surveyed in \citet{nagel2023dynamics} reporting subjective return expectations across individual investors, CFOs, and professional forecasters. 
Table~\ref{tab:nx_regression} summarizes their findings: The first column shows that realized future returns are positively predicted by the dividend-price ratio (a standard return predictor). 
In contrast, past returns do not appear significantly. 
In contrast, regressions for subjective expected returns reveal heterogeneous beliefs: individual investors place a positive and significant weight on past returns (consistent with extrapolation), professional forecasters behave in a contrarian way (a negative coefficient on past returns indicates interpolation), and CFOs lie somewhere in between. 
These results underscore the importance of heterogeneity in expectations among different agents.\footnote{
    See also \citet{atmaz2023contrarians} for evidence of both extrapolative and contrarian behavior. 
    \citet{greenwood_expectations_2014} provides further evidence of extrapolative expectations from past stock return, also associated with the cyclicality of credit and leverage \citep{lopez2017credit}.} 

\begin{table}[t!]
    \centering
    \caption{Subjective return expectations regressions by type of investor.}
    \label{tab:nx_regression}
    \begin{tabular}{lcccc}
   \midrule \midrule
   Dependent var.:         & \textbf{Objective}           & \textbf{Individual}            & \textbf{CFO}            & \textbf{Professional}\\  
   \midrule
   \emph{Predictors}\\
   D/P            & 5.83   &  -0.00  & -0.41  & 0.42\\   
    \hspace{0.2cm}\small{(p-value)} & (0.00) &  (1.00) & (0.31) & (0.73)\\   
   $R_{past}^e$   & 0.36   &  0.87   & 0.30   & -2.78\\   
    \hspace{0.2cm}\small{(p-value)} & (0.71) & (0.02)  & (0.29) & (0.01)\\   
   \midrule \midrule
\end{tabular}
\end{table}


For our paper, a key question is whether expectations correlate with actual trading behavior. 
Existing evidence suggests that yes: \citet{greenwood_expectations_2014} shows that subjective expectations correlate with mutual fund flows, and \citet{GMSU_AER} documents a relationship between beliefs and portfolio choices.\footnote{
    In Section~\ref{sec:belief_taxonomy.sec}, we construct a measure of heterogeneity among professional forecasters and show that this measure of disagreement correlates with stock market turnover, particularly during recessions.}

\paragraph{Fact 3. Expectations correlate with hiring decisions.}

% Previous Table 2 (kept for record; commented out)
% \begin{table}[h!]
% \centering
% \caption{I/B/E/S Expectations and Labor}
% \scalebox{0.9}{
% \begin{tabular}[t]{lcccccc}
%  \tabularnewline \midrule \midrule
%    Dependent Variable: & \multicolumn{3}{c}{$Payroll$} & \multicolumn{3}{c}{Number of workers}\\
%    Model:     & (1) & (2) & (3) & (4) & (5) & (6) \\
%    \midrule
% (Intercept)      & 0.388     & 0.274     & 0.364 & 0.392 & 0.143 & 0.170\\
%                     & (0.247)   & (0.223)   & (0.241) & (0.097) & (0.195) & (0.052)\\
% lag earning growth expectation         & 0.133*** & 0.135*** & 0.136*** & 0.061*** & 0.062*** & 0.067***\\
%                    & (0.043)    & (0.043)   & (0.043)  & (0.014) & (0.014) & (0.014)\\
% 12-month lag return &  & 0.474*** &  &  & 0.535*** & \\
% 		&  & (0.119) & &  & (0.054) &\\
% 12-month return &  &  & 0.234* &  &  & 0.326*** \\
%  			&   &  & (0.125)  &  &  & (0.043) \\
%      \midrule
%      \emph{Fit statistics}\\
% Observations      & 1797  & 1797 & 1797 & 1797  & 1797 & 1797\\
% R2                       & 0.081 & 0.090 & 0.083 & 0.081 & 0.090 & 0.083 \\
% Adjusted R2        & 0.038 & 0.046 & 0.039 & 0.038 & 0.046 & 0.039 \\
%   \midrule \midrule
%    \multicolumn{4}{l}{\emph{Newey-West standard-errors in parentheses (4 lags)}}\\
%    \multicolumn{4}{l}{\emph{Signif. Codes: ***: 0.01, **: 0.05, *: 0.1}}\\
%
% \hline \\[-1.8ex]
% \end{tabular}}
%  \label{tab:payroll}
% \end{table}

\begin{table}[t]
\centering
\caption{I/B/E/S Expectations and Labor}
\label{tab:payroll}
\small
\begin{tabular}{lcccccc}
\toprule\toprule
Dependent Variable: & \multicolumn{3}{c}{Payroll} & \multicolumn{3}{c}{Number of workers} \\
Model: & (1) & (2) & (3) & (4) & (5) & (6) \\
\midrule
(Intercept) & 0.491*** & 0.456*** & 0.467*** & 0.177*** & 0.129*** & 0.158*** \\
 & (0.035) & (0.037) & (0.038) & (0.033) & (0.036) & (0.033) \\
Lag earning growth expectation & 0.230*** & 0.187*** & 0.228*** & 0.277*** & 0.219*** & 0.276*** \\
 & (0.037) & (0.038) & (0.036) & (0.070) & (0.078) & (0.070) \\
12-month lag return &  & 0.326*** &  &  & 0.443*** &  \\
 &  & (0.084) &  &  & (0.101) &  \\
12-month return &  &  & 0.178** &  &  & 0.134 \\
 &  &  & (0.079) &  &  & (0.097) \\
\midrule
\emph{Fit statistics} \\
Observations & 1,746 & 1,746 & 1,746 & 1,746 & 1,746 & 1,746 \\
$R^2$ & 0.036 & 0.044 & 0.039 & 0.036 & 0.047 & 0.037 \\
% Adjusted $R^2$ & 0.035 & 0.043 & 0.038 & 0.036 & 0.046 & 0.036 \\
\midrule\midrule
\multicolumn{7}{l}{\emph{Standard errors clustered at the firm level in parentheses.}} \\
\multicolumn{7}{l}{\emph{Significance: *** $p<0.01$, ** $p<0.05$, * $p<0.10$.}} \\
\bottomrule
\end{tabular}
\end{table}


The previous two facts focus on \emph{subjective expectations} among groups of investors. Toward a business cycle theory, a natural follow-up question is whether these expectations are reflected in \emph{real} economic outcomes. Indeed, \citet{GMS16} shows that CFO expectations predict corporate capital investment, and \citet{armona2019home} finds that expectations about home price appreciations influence construction. In this spirit, we present related evidence: investor expectations also predict firm-level \emph{employment} decisions. To demonstrate this, we use two labor-related measures from Compustat Annual: the realized growth in staff expenses (payroll) and the employee count (workers). These two measures are standardized using the same procedure we apply to our earnings data.

\q{GennaioliStyleRegression}{
We follow the approach of \citet{GMS16}: we regress realized employment growth on the firm-level earnings growth expectations drawn from I/B/E/S. We control for past 12-month stock returns (constructed using individual stock prices from CRSP) and contemporaneous returns, clustering standard errors at the firm level. By including past returns, we aim to capture systematic shocks that might otherwise confound the relationship between expectations and subsequent employment growth.

Table \ref{tab:payroll} shows that earnings growth expectations predict realized employment growth in both total payroll and the number of workers. In light of Facts 1 and 2---namely, that earnings growth expectations exhibit greater volatility than actual outcomes and that CFOs behave differently from individual investors---these results provide reassuring evidence that firm employment decisions correlate with investor expectations.
}

\q{GennaioliCaveat}{
Naturally, this correlation alone does not establish causality: it may simply reflect that the surveyed analysts' information overlaps with managers' private signals, even if investors do not directly influence managers. However, \citet{GMS16} finds that investor expectations continue to affect capital investment even after controlling for managers' survey responses, suggesting a meaningful role for such external perceptions. Unlike \citet{GMS16}, we cannot control for CFO expectations, as we lack sufficient data matches for the employment variables. Motivated by the three facts, we proceed to the model.
}